\section{The literal phase graph as a weighted source and its exact Schur correction}
\label{sec:pgs}

Keep the actual counterfactual full quartet, its multiplicities and its
original primitive from QCV, QPS and MBG:
\[
\begin{gathered}
 m,k\in\mathbb Z_{>0},\quad d=4m,\quad
 \rho=\tfrac12+\delta+i\gamma,\quad0<\delta<\tfrac12,\quad\gamma>2,\\
 h(s)=\prod_{r\in\{\rho,\bar\rho,1-\rho,1-\bar\rho\}}(s-r)^m,
 \quad \Phi'=h,\quad\Phi(0)=0,\quad v_h=g/h,\quad g=2\xi,\\
 \upsilon_h=j_h(v_h),\quad E_h=\mathbb{C}[s]/(h),\quad
 C=\Phi(A_h),\quad C_k=C^{[k]},\quad
 \Phi(\tfrac12+it)=A+iq_\Phi(t).
\end{gathered}
\tag{PGS.1}\label{eq:pgs-original}
\]
Here \(A>0\) is the retained primitive constant, and
\(q_\Phi\) is the original real odd polynomial of degree \(d+1\),
with leading coefficient \(1/(d+1)\). Its derivative is exactly
\(
 [((t-\gamma)^2+\delta^2)((t+\gamma)^2+\delta^2)]^m>0
\).
No replacement of \(h\), its unit or its primitive is made. An
actual off-line zero is not asserted to have been located. The present
calculation evaluates the original source attached to this actual
counterfactual packet, using the full-divisor hypothesis defining that
packet.

The Hilbert space is
\(\mathscr H_k=L^2(\mathbb{R}^k,\prod_i dt_i/(2\pi))\), and
inner products are conjugate linear in the first variable. Put
\[
\begin{gathered}
 s_i=\tfrac12+it_i,\quad\Sigma=\sum_i s_i,\quad
 \tau_k=\sum_i\Phi(s_i)=kA+i\sum_iq_\Phi(t_i),\\
 \mathsf T_k f=\tau_k f,\quad
 \operatorname{dom}\mathsf T_k=\{f:\tau_kf\in\mathscr H_k\},\\
 \mathsf S w=v_h(s)\operatorname{rem}_h(\upsilon_h^{-1}w;s),
 \quad\mathsf S_k=\mathsf S^{\otimes k},\quad
 G_k=\mathsf S_k^*\mathsf S_k .
\end{gathered}
\tag{PGS.2}\label{eq:pgs-source}
\]
All polynomial columns used below are in the indicated multiplication
domains by the original Mellin decay estimates. In particular every
adjoint below is the adjoint of a map with finite dimensional source;
an expression such as \(V^*\mathsf T_k^*L\) denotes its literal
integral pairing, whose columns are in the required domains.

Retain the original full cyclic annihilator and observation:
\[
\begin{gathered}
 \chi=\chi_{h,k},\quad q=\deg\chi=[1+k(m-1)](k+1)^2,
 \quad E=\mathbb{C}[S]/(\chi),\\
 \eta:E\hookrightarrow E_h^{\otimes k},\qquad
 \eta[P]=\upsilon_h^{\otimes k}P(A_h^{[k]})1,\\
 \mathcal P_N=\mathbb{C}[S]_{\le N},\quad N\ge q-1,
 \quad n=N+1,\quad r=N-q+1,\quad\mathcal P_{-1}=0,\\
 B_N:\mathcal P_{N-q}\longrightarrow\mathcal P_N, Q\longmapsto\chi Q,
 \quad J_N:\mathcal P_N\longrightarrow E, P\longmapsto[P],\\
 s_N:E\longrightarrow\mathcal P_N,\quad
 [P]\longmapsto\operatorname{rem}_\chi P.
\end{gathered}
\tag{PGS.3}\label{eq:pgs-sequence}
\]
All finite matrices use increasing monomial frames, including the
original increasing remainder frame of \(E\). Monic division proves
\(\ker J_N=\operatorname{im}B_N\), \(B_N\) injective and
\(J_Ns_N=I_q\). The full length and injectivity of \(\eta\) are
those of MBG.23: each ordered tuple has nilpotence length
\(1+k(m-1)\), its terminal coefficient is
\([k(m-1)]!/((m-1)!)^k\ne0\), and the \((k+1)^2\) arithmetic
centres are distinct. This retains all ordered tuple contributions
before passing to the stated cyclic subspace.

Write \(\mathfrak r\) for the successive remainders modulo the
original \(h(s_i)\). Define the following maps with their full units:
\[
\begin{gathered}
 \mathcal V_N:\mathcal P_N\longrightarrow\mathscr H_k,\quad
 \mathcal V_NP=\Bigl(\prod_i v_h(s_i)\Bigr)P(\Sigma),\\
 L:E\longrightarrow\mathscr H_k,\quad L=\mathsf S_k C_k\eta,
 \qquad L[S^b]=\Bigl(\prod_i v_h(s_i)\Bigr)F_b(s_1,\ldots,s_k),\\
 F_b=\mathfrak r\bigl(\tau_k\Sigma^b\bigr),\quad0\le b<q,
 \qquad\mathsf D_N=\mathsf T_k\mathcal V_N-LJ_N.
\end{gathered}
\tag{PGS.4}\label{eq:pgs-VL}
\]
The formula for \(F_b\) is the exact composite
\(\operatorname{rem}_{\mathfrak h}\upsilon_h^{-\otimes k}
C_k\upsilon_h^{\otimes k}\), proved in MBG.25; every factor
\(v_h(s_i)\) remains in the source. The map \(\mathcal V_N\) is
injective, since a nonzero polynomial in \(\Sigma\) cannot vanish
on a product of open intervals where all \(v_h(s_i)\ne0\).
The map \(L\) is injective as well. Indeed \(\eta\) and
\(\mathsf S_k\) are injective, and the already proved original
critical-value calculation gives
\(\operatorname{Re}c_{k,a,b}\ge
kA[1-(65/289)^m]>0\).
The actual inverse is explicit on the full ordered tensor blocks.
Let \(e_r\) be the original CRT idempotent of \(r\) in \(E_h\),
and for every ordered tuple \(\boldsymbol r=(r_1,\ldots,r_k)\)
put \(e_{\boldsymbol r}=e_{r_1}\otimes\cdots\otimes e_{r_k}\)
and \(c_{\boldsymbol r}=\sum_i\Phi(r_i)\). Since \(\Phi'=h\),
at a root of order \(m\) one has
\(\Phi(s)-\Phi(r)\in(s-r)^{m+1}\), which vanishes in its full
local quotient modulo \((s-r)^m\). Therefore multiplication by
the phase is exactly scalar on every entire primary block, and
\[
 C_k=\sum_{\boldsymbol r}c_{\boldsymbol r}M_{e_{\boldsymbol r}},
 \qquad
 C_k^{-1}=\sum_{\boldsymbol r}c_{\boldsymbol r}^{-1}
                                  M_{e_{\boldsymbol r}}.
\]
Here \(M_e\) is multiplication by the stated original idempotent.
The idempotents are pairwise orthogonal and sum to one, so composing
the two displayed matrices proves their inverse identity directly.
Every ordered tuple is retained, including distinct tuples with the
same critical sum. The arithmetic nilpotents of \(A_h^{[k]}\)
remain those retained in (PGS.3). This proves
the injectivity used here on the full original tensor object.

\subsection{The graph factorization and its finite-rank correction}

Define maps into the same ordered two-component Hilbert space:
\[
\begin{gathered}
 U_N=(\mathcal V_N,\mathsf T_k\mathcal V_N):
     \mathcal P_N\longrightarrow\mathscr H_k\oplus\mathscr H_k,
 \qquad Y=(0,L):E\longrightarrow\mathscr H_k\oplus\mathscr H_k,\\
 \mathcal A_N=(\mathcal V_N,\mathsf D_N)=U_N-YJ_N,\\
 W_N=U_N^*U_N
    =\mathcal V_N^*(I+\mathsf T_k^*\mathsf T_k)\mathcal V_N
       \in\operatorname{Mat}_{n\times n}(\mathbb{C}),\\
 Z_N=U_N^*Y=\mathcal V_N^*\mathsf T_k^*L
       \in\operatorname{Mat}_{n\times q}(\mathbb{C}),\qquad
 K=L^*L=\eta^*C_k^*G_kC_k\eta
       \in\operatorname{Mat}_{q\times q}(\mathbb{C}).
\end{gathered}
\tag{PGS.5}\label{eq:pgs-UYZ}
\]
This \(K\) denotes the displayed cyclic phase-source Gram; it is
not the one-factor covariance matrix called \(K\) in MBG.7.
Both \(W_N\) and \(K\) are positive definite by the injectivities
just proved. Taking the adjoint of the displayed literal graph map
gives, with every cross term and its sign retained,
\[
\boxed{\quad
 \widehat M_N=\mathcal A_N^*\mathcal A_N
     =W_N-Z_NJ_N-J_N^*Z_N^*+J_N^*KJ_N .\quad}
\tag{PGS.6}\label{eq:pgs-gram}
\]
This is precisely \(M_N+\mathsf D_N^*\mathsf D_N\), where
\(M_N=\mathcal V_N^*\mathcal V_N\) is the original AT source
Gram. The correction to \(W_N\) factors as
\[
 \widehat M_N-W_N=
 [\,Z_N\ J_N^*\,]
 \begin{pmatrix}0&-I_q\\-I_q&K\end{pmatrix}
 \begin{bmatrix}Z_N^*\\J_N\end{bmatrix}.
\tag{PGS.7}\label{eq:pgs-rank}
\]
Its rank is at most \(2q\); no assertion that this difference
has a sign is made. The exact relation square is
\[
 \mathcal A_NB_N=U_NB_N,\qquad
 J_NB_N=0,\qquad
 \widehat H_N:=B_N^*\widehat M_NB_N=B_N^*W_NB_N=:H_N^W .
\tag{PGS.8}\label{eq:pgs-relations}
\]
It follows directly by applying the maps to each \(Q\), rather
than discarding the relations. The first components are the same
nonzero original source \(\mathcal V_N(\chi Q)\); the second
components equal \(\mathsf T_k\mathcal V_N(\chi Q)\).
Their original theta primitives and the two cancelling tensor
signs \((-1)^{i-1}(-1)^{i-1}=+1\) are the exact successive
divisions of MBG.26. Thus (PGS.8) is a statement about the original
relation arrows of AT.3 and MBG, including their source norms.

Here every matrix entry is a finite expression in the original moments.
Put \(M^{(h)}_{ab}=\int\bar s^a s^b|v_h(s)|^2dt/(2\pi)\),
\(P_{\alpha j}=[s^\alpha]\Sigma^j\),
\(T_{\alpha j}=[s^\alpha]\tau_k\Sigma^j\), and
\(f_{\alpha b}=[s^\alpha]F_b\). Then
\[
\begin{aligned}
 (W_N)_{ij}
 &=\sum_{\alpha,\beta}
  (\bar P_{\alpha i}P_{\beta j}
       +\bar T_{\alpha i}T_{\beta j})
             \prod_{a=1}^k M^{(h)}_{\alpha_a\beta_a},\\
 (Z_N)_{ib}
 &=\sum_{\alpha,\beta}\bar T_{\alpha i}f_{\beta b}
             \prod_{a=1}^k M^{(h)}_{\alpha_a\beta_a},\\
 K_{bc}&=\sum_{\alpha,\beta}\bar f_{\alpha b}f_{\beta c}
             \prod_{a=1}^k M^{(h)}_{\alpha_a\beta_a}.
\end{aligned}
\tag{PGS.9}\label{eq:pgs-entries}
\]
All sums are finite polynomial expansions. The coefficients
\(f_{\alpha b}\) are obtained by monic division in the original
variables and satisfy \(\alpha_a<d\). Fubini applies because
each integrand is a polynomial times the original product density
with all moments finite. This proves the entry formulas including
the cross source matrix \(Z_N\), rather than replacing it by a norm.

\subsection{The two exact least-norm sections}

For conciseness suppress \(N\) within the next identities. Define
\[
 R_W=s-B(H^W)^{-1}B^*Ws,\qquad
 G_W=R_W^*WR_W=(JW^{-1}J^*)^{-1}.
\tag{PGS.10}\label{eq:pgs-weight-section}
\]
Every lift of \(x\in E\) is \(sx+By\). Completing its positive
\(W\)-quadratic form gives the displayed minimizer, so
\(JR_W=I_q\), \(B^*WR_W=0\). These equalities imply
\(WR_W=J^*G_W\): decompose each vector as a relation plus
\(R_WJp\) and pair against \(R_Wx\). Multiplying by \(W^{-1}\)
and then \(J\) proves the inverse expression in (PGS.10).
The inverse also has the exact decomposition
\[
 W^{-1}=B(H^W)^{-1}B^*+R_WG_W^{-1}R_W^*.
\tag{PGS.11}\label{eq:pgs-inverse-decomposition}
\]
To prove it, use the invertible ordered frame \([R_W\ B]\).
Its \(W\)-Gram is \(\operatorname{diag}(G_W,H^W)\).
Inverting that congruence yields (PGS.11).

For the literal phase graph write a lift as \(p=R_Wx+By\).
Using (PGS.6), (PGS.8) and \(B^*WR_W=0\), its squared norm is
\[
\begin{aligned}
 \|\mathcal A p\|^2
  &=x^*(G_W-R_W^*Z-Z^*R_W+K)x
    +y^*H^Wy-y^*B^*Zx-x^*Z^*By\\
  &=\bigl(y-(H^W)^{-1}B^*Zx\bigr)^*H^W
       \bigl(y-(H^W)^{-1}B^*Zx\bigr)
     +x^*\widehat Gx,
\end{aligned}
\tag{PGS.12}\label{eq:pgs-complete-square}
\]
where consequently the unique graph section and its quotient Gram are
\[
\boxed{\begin{aligned}
 \widehat R&=R_W+B(H^W)^{-1}B^*Z,\\
 \widehat G&=G_W-R_W^*Z-Z^*R_W+K
                    -Z^*B(H^W)^{-1}B^*Z.
\end{aligned}}
\tag{PGS.13}\label{eq:pgs-graph-section}
\]
The signs come from the literal subtraction \(U-YJ\). In particular
\(J\widehat R=I_q\),
\(B^*\widehat M\widehat R=0\) and
\(\widehat G=\widehat R^*\widehat M\widehat R
=(J\widehat M^{-1}J^*)^{-1}\), by the same argument as for
(PGS.10). Thus these are exactly the MBG.30 section and quotient,
not a new quotient chosen for this calculation.

The determinant orientation also stays fixed. In monomial frames
\(\det[s\ B]=1\); replacing \(s\) by either section adds only
relation columns and preserves that determinant. The order \([B\ s]\)
has determinant \((-1)^{qr}\), as in AT.10. Hence
\[
 \det G_W=\frac{\det W}{\det H^W},\qquad
 \det\widehat G=\frac{\det\widehat M}{\det H^W},\qquad
 \frac{\det\widehat G}{\det G_W}
      =\frac{\det\widehat M}{\det W}.
\tag{PGS.14}\label{eq:pgs-determinants}
\]
For \(N=q-1\), the relation domain has dimension zero, its
determinant is one, and every composition through it is zero.
All identities (PGS.8)--(PGS.14) therefore hold there literally.

\subsection{The positive residual and its exact source meaning}

Put
\[
 F_N^{\rm obs}=J_NW_N^{-1}Z_N,\qquad
 \Delta_N=K-Z_N^*W_N^{-1}Z_N,\qquad
 \Pi_N=U_NW_N^{-1}U_N^*.
\tag{PGS.15}\label{eq:pgs-residual-data}
\]
The operator \(\Pi_N\) is the orthogonal projection onto
\(\operatorname{im}U_N\): its adjoint equals itself, its square
equals itself because \(U_N^*U_N=W_N\), and it acts as the
identity on \(U_Np\). It follows that
\[
 \Delta_N=Y^*(I-\Pi_N)Y,\qquad
 x^*\Delta_Nx=\min_{p\in\mathcal P_N}
   \left(\|\mathcal V_Np\|^2+
       \|\mathsf T_k\mathcal V_Np-Lx\|^2\right)>0
       \quad(x\ne0).
\tag{PGS.16}\label{eq:pgs-residual-positive}
\]
The unique unrestricted minimizing coefficient is
\(W_N^{-1}Z_Nx\), obtained by completing this quadratic form.
For strict positivity, equality would put \(Yx=U_Np\); its first
component gives \(\mathcal V_Np=0\), hence \(p=0\), and its
second then gives \(Lx=0\), hence \(x=0\). This proves the strict
sign for every actual finite window.

There is also an explicit positive lower bound independent of \(N\).
Define the full-source form
\[
 K_{\rm low}=L^*(I+\mathsf T_k\mathsf T_k^*)^{-1}L,\qquad
 (K_{\rm low})_{bc}=\int_{\mathbb{R}^k}
   \frac{\overline{F_b(s)}F_c(s)
          \prod_i|v_h(s_i)|^2}
        {1+k^2A^2+(\sum_iq_\Phi(t_i))^2}
                  \prod_i\frac{dt_i}{2\pi}.
\tag{PGS.17}\label{eq:pgs-residual-bound}
\]
For arbitrary \(f\in\operatorname{dom}\mathsf T_k\) and \(y=Lx\),
the pointwise identity is
\[
 |f|^2+|\tau_k f-y|^2
  =(1+|\tau_k|^2)
       \left|f-\frac{\overline{\tau_k}y}{1+|\tau_k|^2}\right|^2
          +\frac{|y|^2}{1+|\tau_k|^2}.
\tag{PGS.18}\label{eq:pgs-pointwise-square}
\]
Its minimizing function lies in the multiplication domain because
\(|\tau_k|/(1+|\tau_k|^2)\le1/2\) and
\(|\tau_k|^2/(1+|\tau_k|^2)\le1\). Integrate and restrict to
\(f=\mathcal V_Np\) to obtain
\[
 0<K_{\rm low}\preceq\Delta_N\preceq K.
\tag{PGS.19}\label{eq:pgs-uniform-N}
\]
The lower matrix is strictly positive because \(L\) is injective
and its scalar denominator is finite and strictly positive almost
everywhere. The upper inequality follows from \(\Pi_N\ge0\).
Thus the residual has a fully specified original source form; it
has not been assigned a bound by arithmetic purity.

Equations (PGS.11) and (PGS.15) turn the quotient into the second
exact completed square
\[
\boxed{\begin{aligned}
 \widehat R_N&=W_N^{-1}Z_N+
                   R_{W,N}(I_q-F_N^{\rm obs}),\\
 \widehat G_N&=(I_q-F_N^{\rm obs})^*G_{W,N}
                   (I_q-F_N^{\rm obs})+\Delta_N.
\end{aligned}}
\tag{PGS.20}\label{eq:pgs-quotient-square}
\]
Indeed \(R_{W,N}^*Z_N=G_{W,N}F_N^{\rm obs}\) and
\(Z_N^*W_N^{-1}Z_N=Z_N^*B_N(H_N^W)^{-1}B_N^*Z_N
 +(F_N^{\rm obs})^*G_{W,N}F_N^{\rm obs}\).
Substitution into (PGS.13) proves both lines with their matrix order.

For the actual commuting source diagram, define
\(\mathcal E_N=(I-\Pi_N)Y\). Then
\[
 \mathcal A_N\widehat R_N
       =U_NR_{W,N}(I_q-F_N^{\rm obs})-\mathcal E_N,\qquad
 (U_NR_{W,N})^*\mathcal E_N=0,\qquad
 \mathcal E_N^*\mathcal E_N=\Delta_N.
\tag{PGS.21}\label{eq:pgs-orthogonal-source-map}
\]
Use \(U_NW_N^{-1}Z_N=\Pi_NY\) and \(J_N\widehat R_N=I\)
to prove its first identity; the remaining two are projection
identities. This is an exact linear map of the cyclic quotient into
two specified orthogonal parts of the original graph Hilbert target.
It proves the morphism underlying (PGS.20) even when
\(I_q-F_N^{\rm obs}\) is singular. No invertibility of that matrix
has been presumed or needed.

All coefficient maps just displayed extend to each original active
mask by applying the same linear map coordinate by coordinate. For
example \(\ker J_N=(\chi\mathcal P_{N-q})\),
\(\ker B_N=\ker\widehat R_N=0\), and
\(\ker\mathcal E_N=0\) by \(\Delta_N>0\). For
\(F_N^{\rm obs}\) the supported-zero fibre is its literal kernel;
it is not asserted injective. Every present empty-mask tuple stays
present and maps to the empty tuple; \(\tau_{\rm out}\) maps to
\(\tau_{\rm out}\). At an original compatible realization
\(R\to\mathbb{C}\), these are \(R\)-linear maps and commute with
scalar extension on pure tensors. The realization and its full kernel
are exactly MBG.34a--34b. Thus none of the new norm or section maps
identifies a supported zero with absence or adds a residue-field map
into \(\mathbb{C}\).

\subsection{Literal transitions between polynomial windows}

For \(M\ge N\), let \(e_{N,M}:\mathcal P_N\hookrightarrow
\mathcal P_M\) be the original coefficient inclusion, and let
\(e_{N-q,M-q}\) be the inclusion of relation coefficient spaces.
Polynomial evaluation proves all four commuting identities
\[
\begin{gathered}
 U_Me_{N,M}=U_N,\quad
 J_Me_{N,M}=J_N,\quad
 B_Me_{N-q,M-q}=e_{N,M}B_N,\quad
 \mathcal A_Me_{N,M}=\mathcal A_N,\\
 W_N=e_{N,M}^*W_Me_{N,M},\qquad
 Z_N=e_{N,M}^*Z_M,\qquad
 \widehat M_N=e_{N,M}^*\widehat M_Me_{N,M}.
\end{gathered}
\tag{PGS.22}\label{eq:pgs-window-squares}
\]
In particular \(\Pi_N\preceq\Pi_M\), since their ranges are
nested. Orthogonal projection onto the range difference gives the
exact residual drop
\[
 \Delta_N-\Delta_M=Y^*(\Pi_M-\Pi_N)Y\succeq0.
\tag{PGS.23}\label{eq:pgs-residual-transition}
\]
For the quotient sections the full discrepancy, including its
relation source, is
\[
\begin{gathered}
 T_{N,M}=(H_M^W)^{-1}B_M^*\widehat M_Me_{N,M}\widehat R_N,\\
 e_{N,M}\widehat R_N-\widehat R_M=B_MT_{N,M},\qquad
 \widehat G_N-\widehat G_M=T_{N,M}^*H_M^WT_{N,M}\succeq0.
\end{gathered}
\tag{PGS.24}\label{eq:pgs-section-transition}
\]
The first right side is the relation-orthogonal projection of the
given lift onto \(\operatorname{im}B_M\); subtracting it yields
the unique least-norm lift in the larger window. Its graph is
orthogonal to the relation graph, so Pythagoras proves the last
identity. The same identities with \(\widehat M,\widehat R,
\widehat G\) replaced by \(W,R_W,G_W\) follow from the same
projection, with no changes to \(B_M\), \(J_M\), or their labels.

Here is an explicit rank-one version of every consecutive transition.
In \(\mathcal P_{N+1}\) set
\[
\begin{gathered}
 b_N=\chi S^{N+1-q},\qquad
 d_N=b_N-e_{N,N+1}B_N(H_N^W)^{-1}
                 B_N^*e_{N,N+1}^*W_{N+1}b_N,\\
 a_N=d_N^*W_{N+1}d_N>0,\qquad
 \beta_N=d_N^*W_{N+1}e_{N,N+1}R_{W,N},\\
 \alpha_N=\beta_N-d_N^*Z_{N+1}.
\end{gathered}
\tag{PGS.25}\label{eq:pgs-step-data}
\]
Here \(d_N\) is a column of length \(N+2\), and
\(\alpha_N,\beta_N\) are rows of length \(q\). The polynomial
\(d_N\) has leading coefficient one in degree \(N+1\), so it
is nonzero, and it is \(W_{N+1}\)-orthogonal to all old relation
columns. New and old relation graphs are consequently orthogonal
by (PGS.8). Pairing \(U_{N+1}d_N\) against
\(\mathcal A_N\widehat R_N\) gives exactly \(\alpha_N\):
the term \(B_N(H_N^W)^{-1}B_N^*Z_N\) in (PGS.13) pairs to
zero, leaving \(\beta_N-d_N^*Z_{N+1}\).
Subtracting this single projection proves
\[
\begin{aligned}
 \widehat R_{N+1}
    &=e_{N,N+1}\widehat R_N-d_Na_N^{-1}\alpha_N,\qquad
 &\widehat G_{N+1}
    &=\widehat G_N-a_N^{-1}\alpha_N^*\alpha_N,\\
 R_{W,N+1}
    &=e_{N,N+1}R_{W,N}-d_Na_N^{-1}\beta_N,\qquad
 &G_{W,N+1}
    &=G_{W,N}-a_N^{-1}\beta_N^*\beta_N.
\end{aligned}
\tag{PGS.26}\label{eq:pgs-rank-one}
\]
Both proposed section corrections are relations since
\(J_{N+1}d_N=0\). They kill exactly the remaining new relation
pairing, and all old pairings already vanish. This proves the
section and Gram formulas without an assumed canonical transport.

\subsection{All four signed endpoints, with the finite correction evaluated}

Set
\[
 \widehat\sigma_N=
    a_N^{-1}\alpha_N\widehat G_N^{-1}\alpha_N^*,\qquad
 \sigma_N^W=a_N^{-1}\beta_NG_{W,N}^{-1}\beta_N^*.
\tag{PGS.27}\label{eq:pgs-sigmas}
\]
These are real numbers in \([0,1)\). Nonnegativity is the
positive inverse Gram. The actual congruent matrix
\(\widehat G_N^{-1/2}\widehat G_{N+1}\widehat G_N^{-1/2}\)
equals \(I_q-a_N^{-1}(\widehat G_N^{-1/2}\alpha_N^*)
(\alpha_N\widehat G_N^{-1/2})\), and the weighted matrix has
the same expression with \(\alpha_N,\widehat G_N\) replaced by
\(\beta_N,G_{W,N}\). These congruent matrices have eigenvalues
one except possibly \(1-\widehat\sigma_N\), or
\(1-\sigma_N^W\); their strict positivity proves the strict upper
bound. The rank-one determinant identity, proved by extending its
nonzero update vector to a unitary basis, gives
\[
 \frac{\det\widehat G_{N+1}}{\det\widehat G_N}
       =1-\widehat\sigma_N,\qquad
 \frac{\det G_{W,N+1}}{\det G_{W,N}}
       =1-\sigma_N^W.
\tag{PGS.28}\label{eq:pgs-step-determinants}
\]
If the update vector is zero both equalities read one, so the proof
includes that case.

For any of the compatible source families \(H\), retain exactly
the original AT.14 convention
\[
 \mathcal B(H)=\log\det G^H_{q-1}+\log\det G^H_q
                 -\log\det G^H_{2q-1}-\log\det G^H_{2q}.
\tag{PGS.29}\label{eq:pgs-four-signs}
\]
Define weights \(\omega_{q-1}=\omega_{2q-1}=1\), and
\(\omega_j=2\) for \(q\le j\le2q-2\); for \(q=1\)
this interior is empty. Telescoping (PGS.28) along the two
intervals \([q-1,2q-1]\) and \([q,2q]\) proves the literal
finite evaluation
\[
\begin{aligned}
 \mathcal B(\widehat M)
       &=-\sum_{j=q-1}^{2q-1}\omega_j
                          \log(1-\widehat\sigma_j)\ \ge0,\\
 \mathcal B(W)
       &=-\sum_{j=q-1}^{2q-1}\omega_j
                          \log(1-\sigma_j^W)\ \ge0,\\
 \mathcal B(\widehat M)-\mathcal B(W)
       &=\sum_{j=q-1}^{2q-1}\omega_j
                  \log\frac{1-\sigma_j^W}{1-\widehat\sigma_j}.
\end{aligned}
\tag{PGS.30}\label{eq:pgs-four-rank-one}
\]
Thus its inputs are the explicit moment matrices (PGS.9), the
weighted monic relation columns (PGS.25), and the computed cross
rows \(\alpha_j=\beta_j-d_j^*Z_{j+1}\). No sign of their
difference is inferred from the separate nonnegative expressions.
For \(q=1\), (PGS.29) cancels its two occurrences of endpoint
one, and (PGS.30) is exactly the sum of the steps zero-to-one and
one-to-two. The same formulas handle the original quartet, for
which \(q\) takes its full value in (PGS.3).

Equivalently, the quotient-sized matrices
\[
 X_N=G_{W,N}^{-1/2}
       \bigl((I-F_N^{\rm obs})^*G_{W,N}(I-F_N^{\rm obs})
                           +\Delta_N\bigr)G_{W,N}^{-1/2}>0
\tag{PGS.31}\label{eq:pgs-quotient-ratio}
\]
give the same answer as
\[
\begin{aligned}
 \mathcal B(\widehat M)-\mathcal B(W)
   &=\log\det X_{q-1}+\log\det X_q
        -\log\det X_{2q-1}-\log\det X_{2q},\\
 \det X_N&=\det\widehat M_N/\det W_N.
\end{aligned}
\tag{PGS.32}\label{eq:pgs-four-ratio}
\]
These follow from (PGS.14) and (PGS.20); the second is also
the exact finite-rank determinant correction in (PGS.7).

The original Gamma source is still
\[
 m_k^\Gamma(t)=\frac{(\sqrt{2\pi})^k}{c_{k/4}}
                 \frac{|\Gamma(k/4+it/2)|^2}{2\pi},\qquad
 c_a=2^{1-2a}\Gamma(2a),\qquad S=k/2+it.
\tag{PGS.33}\label{eq:pgs-reference}
\]
Its mass is \((2\pi)^{k/2}\), and the arithmetic mass remains
\(\mu_h^k\), \(\mu_h=\int |v_h(1/2+it)|^2dt/(2\pi)\).
With exactly those families, the full three-part arithmetic identity is
\[
\begin{aligned}
 \mathcal B(\widehat M)-\mathcal B(M^\Gamma)
  ={}&\mathcal B(M)-\mathcal B(M^\Gamma)\\
    &+\mathcal B(W)-\mathcal B(M)\\
    &+\sum_{j=q-1}^{2q-1}\omega_j
             \log\frac{1-\sigma_j^W}{1-\widehat\sigma_j}.
\end{aligned}
\tag{PGS.34}\label{eq:pgs-three-part}
\]
The first part is the original AT arithmetic/reference transport.
The second is evaluated from the original weighted scalar density in
the accompanying phase-density calculation. The final part is the
literal phase observation correction calculated here. This equality
retains all relation determinants, full units, reference masses,
ordered endpoints and coefficient maps.

\subsection{The exact scalar source, its fibre residual and all cross convolutions}

The scalar reduction has an explicit isometry on the full Hilbert
source, including the phase observation. Absorb only the original
measure factor into the notation
\(w(t)=|v_h(1/2+it)|^2/(2\pi)\); this changes no source vector.
For an integrable function on the product line write its fibre integral
over \(t_1+\cdots+t_k=t\) by substituting
\(t_k=t-\sum_{i<k}t_i\) and integrating in
\(dt_1\cdots dt_{k-1}\). At \(k=1\) this means evaluation at \(t\).
Put \(m_k=w^{*k}\). Expanding the actual scalar multiplier gives
\[
\begin{aligned}
 r_k(t)&=\int_{\sum t_i=t}
      \bigl(1+k^2A^2+(\sum_iq_\Phi(t_i))^2\bigr)
                                  \prod_i w(t_i),\\
 r_1(t)&=(1+A^2+q_\Phi(t)^2)w(t),\\
 r_k&=(1+k^2A^2)w^{*k}
       +k(q_\Phi^2w)*w^{*(k-1)}\\
     &\hspace{17mm}
       +k(k-1)(q_\Phi w)*(q_\Phi w)*w^{*(k-2)}\qquad(k\ge2).
\end{aligned}
\tag{PGS.35}\label{eq:pgs-scalar-density}
\]
Here \(w^{*0}=\delta_0\) is the convolution identity measure,
of mass one; thus the displayed last term at \(k=2\) is exactly
\(2(q_\Phi w)*(q_\Phi w)\).
The \(k\) diagonal terms and \(k(k-1)\) ordered distinct pairs
in the square give exactly these coefficients by Fubini and
permutation of the original tensor variables. The last convolution
is a signed function, whereas the fibre formula proves
\(r_k\ge(1+k^2A^2)m_k>0\) almost everywhere. Every term is
absolutely integrable, including every fixed polynomial moment,
by the original one-variable moments and the triangle inequality.
Both densities are finite almost everywhere, which suffices for
the following specified Hilbert spaces and maps.

The map
\[
\begin{gathered}
 \mathscr J_k:L^2(\mathbb{R},r_k(t)dt)
               \longrightarrow\mathscr H_k\oplus\mathscr H_k,\\
 \mathscr J_k f=
   \left(\Bigl(\prod_i v_h(s_i)\Bigr)f(\sum_i t_i),
    \tau_k\Bigl(\prod_i v_h(s_i)\Bigr)f(\sum_i t_i)\right),\\
 \mathscr J_k^*\mathscr J_k=I,\qquad
 U_N=\mathscr J_k\iota_N,\quad
 (\iota_NP)(t)=P(k/2+it).
\end{gathered}
\tag{PGS.36}\label{eq:pgs-scalar-isometry}
\]
is well-defined on equivalence classes: a null set for \(r_kdt\)
has zero product integral of both displayed squared components,
by (PGS.35). The same formula gives equality of squared norms;
polarization gives equality of pairings. Its range is therefore
closed and its stated adjoint identity holds. In particular
\(W_N=\iota_N^*\iota_N\) in this exact weighted scalar source.

The full cyclic observation has explicit fibre data too. Define
the row \(z(t)\) and the positive semidefinite matrix
\(\Lambda(t)\) by
\[
 z_b(t)=\int_{\sum t_i=t}\overline{\tau_k}F_b(s)\prod_iw(t_i),
 \qquad
 \Lambda_{bc}(t)=\int_{\sum t_i=t}\overline{F_b(s)}F_c(s)
                                                    \prod_iw(t_i).
\tag{PGS.37}\label{eq:pgs-fibre-data}
\]
These are concrete finite sums of one-variable convolutions.
For example, with the same \(f_{\alpha b}\) as in (PGS.9), set
\(u_a(t)=s(t)^aw(t)\),
\(b_a(t)=(A-iq_\Phi(t))s(t)^aw(t)\), and
\(u_{a,c}(t)=\bar s(t)^as(t)^cw(t)\). Then
\[
\begin{aligned}
 z_b&=\sum_{\alpha}f_{\alpha b}\sum_{i=1}^k
       u_{\alpha_1}*\cdots*u_{\alpha_{i-1}}*b_{\alpha_i}
               *u_{\alpha_{i+1}}*\cdots*u_{\alpha_k},\\
 \Lambda_{bc}&=\sum_{\alpha,\beta}\bar f_{\alpha b}f_{\beta c}
                u_{\alpha_1,\beta_1}*\cdots*u_{\alpha_k,\beta_k},\\
 (Z_N)_{ab}&=\int\overline{(k/2+it)^a}z_b(t)\,dt,\qquad
 K=\int\Lambda(t)\,dt.
\end{aligned}
\tag{PGS.38}\label{eq:pgs-cross-convolutions}
\]
At \(k=1\) each convolution list is its one factor. Expansion
of the original polynomials, followed by the defining fibre
integral, proves every identity. No phase source coordinate is
replaced by an independently chosen scalar function.

For \(x\in E\), fibre Cauchy--Schwarz applied to
\(\tau_k\) and \(\sum_bF_bx_b\) gives
\[
 |z(t)x|^2\le(r_k(t)-m_k(t))x^*\Lambda(t)x.
\tag{PGS.39}\label{eq:pgs-fibre-CS}
\]
Indeed the first squared factor integrates to
\(\int_{\sum t_i=t}|\tau_k|^2\prod_iw=r_k-m_k\),
and the second to \(x^*\Lambda x\). Thus
\(g_x(t)=z(t)x/r_k(t)\) is a specified element of \(L^2(r_kdt)\),
setting it to zero on the negligible exceptional set where division
is not defined. Pair against any \(f\) in this scalar Hilbert
space, first on bounded truncations and then by Cauchy--Schwarz:
\(\langle\mathscr J_kf,Yx\rangle=\int\bar f zx
=\langle f,g_x\rangle_{r_k}\). It follows that
\(g=\mathscr J_k^*Y\), as a linear map from \(E\).

Let \(\mathsf P_N=\iota_NW_N^{-1}\iota_N^*\) be the orthogonal
projection onto the actual scalar polynomial space and define
\(Y_\perp=(I-\mathscr J_k\mathscr J_k^*)Y\). The complete
original observation and the finite-window residual now decompose as
\[
\begin{gathered}
 Y=\mathscr J_kg+Y_\perp,\qquad
 \mathscr J_k^*Y_\perp=0,\qquad
 \Omega_k=Y_\perp^*Y_\perp
      =K-\int\frac{z(t)^*z(t)}{r_k(t)}\,dt,\\
 \Delta_N=\Omega_k+
       ((I-\mathsf P_N)g)^*((I-\mathsf P_N)g),\\
 0<\int\frac{m_k(t)}{r_k(t)}\Lambda(t)\,dt
                  \preceq\Omega_k\preceq\Delta_N.
\end{gathered}
\tag{PGS.40}\label{eq:pgs-fibre-residual}
\]
All adjoints in the second line use the original scalar density
\(r_kdt\). The first line is an orthogonal projection identity.
Since \(U_N=\mathscr J_k\iota_N\), its projection is
\(\Pi_N=\mathscr J_k\mathsf P_N\mathscr J_k^*\); substituting
the first decomposition into (PGS.16) proves the second line.
Inequality (PGS.39), divided by \(r_k\) and integrated, proves
the last lower bound. It is strict because \(m_k/r_k>0\)
almost everywhere and
\(\int x^*\Lambda x=x^*Kx>0\) for every nonzero \(x\).
This gives the full map and positive fibre contribution before any
limit of polynomial windows is considered. It does not assume that
the finite observation \(F_N^{\rm obs}\) is an identity.

In particular the exact finite minimizer in (PGS.16) is the scalar
projection \(\mathsf P_Ng_x\), represented in the original
coefficient frame by \(W_N^{-1}Z_Nx\). Combining it with the
specified jet correction \(R_{W,N}(I-F_N^{\rm obs})x\) gives
\(\widehat R_Nx\) exactly as in (PGS.20). Equations
(PGS.35)--(PGS.40) therefore construct the scalar weighted source,
the full mixed observation, their projection, every fibre residual
and their return to the original quotient in a single commuting
diagram of explicitly given maps.

The two proved positive source bounds are themselves related by an
exact morphism and norm equality. For \(x\in E\) write
\(F_x=\sum_bF_bx_b\), \(u=\sum_it_i\), and
\(g_x(u)=z(u)x/r_k(u)\). The scalar-source minimizing vector is
\(\prod_i v_h(s_i)g_x(u)\), whereas the whole multiplication-graph
minimizer from (PGS.18) is
\(\prod_i v_h(s_i)\overline{\tau_k}F_x(s)/(1+|\tau_k|^2)\).
Both belong to the same original multiplication domain. Substituting
the first into (PGS.18) and using (PGS.40) gives
\[
\begin{aligned}
 x^*(\Omega_k-K_{\rm low})x
   =\int_{\mathbb{R}^k}(1+|\tau_k|^2)
      \left|g_x(\sum_it_i)-
        \frac{\overline{\tau_k}F_x(s)}{1+|\tau_k|^2}\right|^2
                                      \prod_iw(t_i)\,dt_i\ \ge0,\\
 0<K_{\rm low}\preceq\Omega_k\preceq\Delta_N,\qquad
 \Omega_1=K_{\rm low}\quad(k=1).
\end{aligned}
\tag{PGS.41}\label{eq:pgs-two-residuals}
\]
Indeed the squared distance from \(Yx\) to the closed scalar graph
is \(x^*\Omega_kx\); the minimum over the entire multiplication
graph is \(x^*K_{\rm low}x\). The pointwise completed square proves
their displayed difference, including the unchanged product measure.
At \(k=1\), (PGS.37) and (PGS.35) give
\(g_x(t)=\overline{\Phi(1/2+it)}F_x(1/2+it)
 /(1+|\Phi(1/2+it)|^2)\) almost everywhere, so the integrand
vanishes. This proves the last equality without identifying the two
minimizers for higher tensor powers.
